\documentclass{article}
\usepackage{amsmath,amssymb,amsthm}
\usepackage{enumitem}
\usepackage{hyperref}
\newtheorem{theorem}{Theorem}
\newtheorem{lemma}{Lemma}
\newtheorem{assumption}{Assumption}
\newcommand{\norm}[1]{\left\lVert #1\right\rVert}
\newcommand{\inner}[2]{\left\langle #1,#2\right\rangle}
\begin{document}

%%%%%%%%%%%%%%%%%%%%%%%%%%%%%%%%%%%%%%%%%%%%%%%%%%%%%%%%%%%%%%%%%%%%%%%%%%%%
\section*{Primal–Dual Hybrid Gradient (PDHG)}
%%%%%%%%%%%%%%%%%%%%%%%%%%%%%%%%%%%%%%%%%%%%%%%%%%%%%%%%%%%%%%%%%%%%%%%%%%%%

\section*{Mathematical Formulation}
We consider the convex–concave saddle–point problem  

\[
\min_{x\in\mathbb R^{n}}\max_{y\in\mathbb R^{m}} \; \mathcal L (x,y)
:= f(x)+\langle Kx,y\rangle - g^{*}(y),
\tag{1}
\]

where  

\begin{itemize}[noitemsep]
\item $f:\mathbb R^{n}\to\mathbb R\cup\{+\infty\}$ is proper, closed and convex;
\item $g^{*}:\mathbb R^{m}\to\mathbb R\cup\{+\infty\}$ is the convex conjugate of a convex $g$;
\item $K\in\mathbb R^{m\times n}$ is a linear operator.
\end{itemize}

Let $\tau>0$ and $\sigma>0$ be primal and dual step sizes satisfying  

\[
\tau\sigma\|K\|^{2}<1.
\tag{2}
\]

Define the extrapolated primal iterate  

\[
\bar x_{k}=x_{k}+\theta (x_{k}-x_{k-1}),\qquad \theta\in[0,1].
\tag{3}
\]

The PDHG updates are  

\[
\boxed{
\begin{aligned}
y_{k+1}&=\operatorname{prox}_{\sigma g^{*}}\!\Bigl(y_{k}+\sigma K\bar x_{k}\Bigr),\\[4pt]
x_{k+1}&=\operatorname{prox}_{\tau f}\!\Bigl(x_{k}-\tau K^{\!\top}y_{k+1}\Bigr),\\[4pt]
\bar x_{k+1}&=x_{k+1}+\theta (x_{k+1}-x_{k}),
\end{aligned}}
\tag{4}
\]

where for a convex $\phi$  

\[
\operatorname{prox}_{\lambda\phi}(v)=\arg\min_{u}\Bigl\{\phi(u)+\frac{1}{2\lambda}\|u-v\|^{2}\Bigr\}.
\]

When $g^{*}=0$ (i.e. $g\equiv0$) the dual step reduces to a simple gradient ascent
$y_{k+1}=y_{k}+\sigma K\bar x_{k}$.

--------------------------------------------------------------------
\section*{Pseudocode}
\begin{verbatim}
Input:   f, g*, K, step sizes τ, σ, extrapolation θ∈[0,1],
         maxIter, tolerance ε
Initialize: x0 ← 0, y0 ← 0, x_{-1} ← x0
for k = 0,1,…,maxIter-1 do
    # extrapolation
    if k==0 then  barx ← x0
    else          barx ← xk + θ*(xk - x_{k-1})
    end
    # dual ascent (prox of g*)
    y_{k+1} ← prox_{σ g*}( y_k + σ * K * barx )
    # primal descent (prox of f)
    x_{k+1} ← prox_{τ f}( x_k - τ * K^T * y_{k+1} )
    # stopping test
    if  norm(x_{k+1}-x_k) < ε  then break
    end
end
return x_{k+1}, y_{k+1}
\end{verbatim}
--------------------------------------------------------------------
\section*{Dry Run Example}
We illustrate (4) on a scalar problem with  

\[
f(x)=(x-3)^{2},\qquad g^{*}\equiv0,\qquad K=1,
\]

so that $\mathcal L(x,y)= (x-3)^{2}+yx$.  
The proximal operator of $f$ with step $\tau$ admits a closed form:

\[
\operatorname{prox}_{\tau f}(v)=
\arg\min_{x}\Bigl\{(x-3)^{2}+ \tfrac{1}{2\tau}(x-v)^{2}\Bigr\}
=\frac{v+2\tau\cdot 3}{1+2\tau}.
\tag{5}
\]

Choose $\tau=\sigma=0.1$, $\theta=1$, and initialise $x_{0}=0$, $y_{0}=0$.
The table below shows three iterations.

\begin{center}
\begin{tabular}{c|c|c|c|c}
$k$ & $y_{k+1}=y_{k}+\sigma \bar x_{k}$ &
$x_{k+1}= \operatorname{prox}_{\tau f}(x_{k}-\tau y_{k+1})$ &
$f(x_{k+1})$ & $\bar x_{k+1}=x_{k+1}+ \theta (x_{k+1}-x_{k})$\\ \hline
0 & $0+0.1\cdot0=0$ &
$\frac{0+0.6}{1.2}=0.5$ &
$(0.5-3)^{2}=6.25$ &
$0.5+1\cdot(0.5-0)=1.0$\\ \hline
1 & $0+0.1\cdot1.0=0.05$ &
$\frac{0.5-0.1\cdot0.05+0.6}{1.2}=0.9125$ &
$(0.9125-3)^{2}\approx4.357$ &
$0.9125+1\cdot(0.9125-0.5)=1.325$\\ \hline
2 & $0.05+0.1\cdot1.325=0.1825$ &
$\frac{0.9125-0.1\cdot0.1825+0.6}{1.2}=1.2486$ &
$(1.2486-3)^{2}\approx3.067$ &
$1.2486+1\cdot(1.2486-0.9125)=1.5847$
\end{tabular}
\end{center}
The objective values decrease, illustrating the contraction guaranteed by the theory below.

%%%%%%%%%%%%%%%%%%%%%%%%%%%%%%%%%%%%%%%%%%%%%%%%%%%%%%%%%%%%%%%%%%%%%%%%%%%%
\section*{Assumptions}
%%%%%%%%%%%%%%%%%%%%%%%%%%%%%%%%%%%%%%%%%%%%%%%%%%%%%%%%%%%%%%%%%%%%%%%%%%%%
\begin{assumption}[Convexity and Smoothness of $f$]\label{ass:f}
$f$ is convex and $L_{f}$‑smooth, i.e.  

\[
\|\nabla f(x)-\nabla f(u)\|\le L_{f}\|x-u\|,\qquad\forall x,u.
\tag{6}
\]
\end{assumption}

\begin{assumption}[Convexity of $g^{*}$]\label{ass:g}
$g^{*}$ is proper, closed and convex. Its proximal operator is assumed to be
computable exactly.
\end{assumption}

\begin{assumption}[Linear Operator]\label{ass:K}
$K\in\mathbb R^{m\times n}$ with operator norm $\|K\|$. The step sizes satisfy  

\[
\tau\sigma\|K\|^{2}<1 .
\tag{7}
\]

\end{assumption}

\begin{assumption}[Bounded Saddle Point]\label{ass:sol}
There exists a saddle point $(x^{\star},y^{\star})$ such that  

\[
\mathcal L(x^{\star},y)\le \mathcal L(x^{\star},y^{\star})\le \mathcal L(x,y^{\star}),
\qquad\forall x,y .
\tag{8}
\]

\end{assumption}

%%%%%%%%%%%%%%%%%%%%%%%%%%%%%%%%%%%%%%%%%%%%%%%%%%%%%%%%%%%%%%%%%%%%%%%%%%%%
\section*{Main Convergence Theorem}
%%%%%%%%%%%%%%%%%%%%%%%%%%%%%%%%%%%%%%%%%%%%%%%%%%%%%%%%%%%%%%%%%%%%%%%%%%%%
\begin{theorem}[Ergodic $O(1/T)$ convergence]\label{thm:main}
Let $\{(x_{k},y_{k})\}_{k\ge0}$ be generated by \eqref{4} with $\theta=1$ and
step sizes satisfying \eqref{7}. Define the ergodic averages  

\[
\bar x_{T}:=\frac{1}{T}\sum_{k=1}^{T}x_{k},\qquad 
\bar y_{T}:=\frac{1}{T}\sum_{k=1}^{T}y_{k}.
\]

Then for any saddle point $(x^{\star},y^{\star})$  

\[
\mathcal G_{T}:=
\sup_{y}\bigl\langle K\bar x_{T},y-y^{\star}\bigr\rangle
   +f(\bar x_{T})-f(x^{\star})
   -\bigl\langle Kx^{\star},\bar y_{T}-y^{\star}\bigr\rangle
   +g^{*}(\bar y_{T})-g^{*}(y^{\star})
\le
\frac{C}{T},
\tag{9}
\]

where  

\[
C=
\frac{1}{2\tau}\|x_{0}-x^{\star}\|^{2}
+\frac{1}{2\sigma}\|y_{0}-y^{\star}\|^{2}.
\tag{10}
\]

Consequently $\mathcal G_{T}\to0$ at the rate $O(1/T)$.
\end{theorem}

\begin{theorem}[Non‑ergodic $O(1/k)$ bound]\label{thm:nonerg}
If additionally $f$ is $L_{f}$‑smooth and $K$ is full‑rank, then  

\[
\|x_{k}-x^{\star}\|^{2}+\|y_{k}-y^{\star}\|^{2}\le\frac{C'}{k},
\qquad C'>0,
\tag{11}
\]

i.e. the iterates themselves converge with the same order.
\end{theorem}

%%%%%%%%%%%%%%%%%%%%%%%%%%%%%%%%%%%%%%%%%%%%%%%%%%%%%%%%%%%%%%%%%%%%%%%%%%%%
\section*{Theoretical Properties}
%%%%%%%%%%%%%%%%%%%%%%%%%%%%%%%%%%%%%%%%%%%%%%%%%%%%%%%%%%%%%%%%%%%%%%%%%%%%
\begin{itemize}
\item \textbf{Stability.} Condition \eqref{7} is necessary and sufficient for the
Lyapunov function (defined below) to be non‑increasing.
\item \textbf{Hyper‑parameter sensitivity.}
The bound $C$ in \eqref{10} scales inversely with $\tau$ and $\sigma$;
larger steps give a smaller $C$ but must respect \eqref{7}.
\item \textbf{Comparison.}
When $g^{*}=0$, PDHG reduces to a primal gradient descent with step
$\tau$ and a dual ascent with step $\sigma$. The $O(1/T)$ rate matches that of
deterministic gradient methods for convex smooth problems (see Example 1).
\end{itemize}

%%%%%%%%%%%%%%%%%%%%%%%%%%%%%%%%%%%%%%%%%%%%%%%%%%%%%%%%%%%%%%%%%%%%%%%%%%%%
\section*{Preliminaries (Definitions and Lemmas)}
%%%%%%%%%%%%%%%%%%%%%%%%%%%%%%%%%%%%%%%%%%%%%%%%%%%%%%%%%%%%%%%%%%%%%%%%%%%%
\begin{lemma}[Three‑point identity for prox]\label{lem:prox}
For any convex $\phi$ and $\lambda>0$,
\[
\inner{u-\operatorname{prox}_{\lambda\phi}(u)}{z-\operatorname{prox}_{\lambda\phi}(u)}
\ge \lambda\bigl(\phi(z)-\phi(\operatorname{prox}_{\lambda\phi}(u))\bigr),
\qquad\forall z.
\tag{12}
\]
\end{lemma}
\begin{proof}
Standard property of proximal operators (see e.g. \cite{parikh2014prox}).
\end{proof}

\begin{lemma}[Descent lemma for $L_{f}$‑smooth functions]\label{lem:descent}
If $f$ satisfies \eqref{6}, then for any $x,u$,
\[
f(u)\le f(x)+\inner{\nabla f(x)}{u-x}
+\frac{L_{f}}{2}\|u-x\|^{2}.
\tag{13}
\]
\end{lemma}
\begin{proof}
Direct integration of the Lipschitz condition on $\nabla f$.
\end{proof}

%%%%%%%%%%%%%%%%%%%%%%%%%%%%%%%%%%%%%%%%%%%%%%%%%%%%%%%%%%%%%%%%%%%%%%%%%%%%
\section*{Main Proof (Step‑by‑Step Derivation)}
%%%%%%%%%%%%%%%%%%%%%%%%%%%%%%%%%%%%%%%%%%%%%%%%%%%%%%%%%%%%%%%%%%%%%%%%%%%%
We prove Theorem~\ref{thm:main}.  Throughout the proof we fix an arbitrary
saddle point $(x^{\star},y^{\star})$.

\noindent\textbf{Step 1.  Lyapunov function.}  
Define  

\[
\Phi_{k}:=
\frac{1}{2\tau}\|x_{k}-x^{\star}\|^{2}
+\frac{1}{2\sigma}\|y_{k}-y^{\star}\|^{2}.
\tag{14}
\]

Our goal is to show $\Phi_{k+1}\le \Phi_{k}-\tau\Delta_{k}$ where
$\Delta_{k}$ contains the primal–dual gap terms.

\noindent\textbf{Step 2.  Expand $\|x_{k+1}-x^{\star}\|^{2}$.}  
Using the update $x_{k+1}= \operatorname{prox}_{\tau f}(x_{k}-\tau K^{\!\top}y_{k+1})$ and
Lemma~\ref{lem:prox} with $u=x_{k}-\tau K^{\!\top}y_{k+1}$ and $z=x^{\star}$ we obtain  

\[
\inner{x_{k+1}-\bigl(x_{k}-\tau K^{\!\top}y_{k+1}\bigr)}{x^{\star}-x_{k+1}}
\ge \tau\bigl(f(x^{\star})-f(x_{k+1})\bigr).
\tag{15}
\]

Rearranging \eqref{15} yields  

\[
\|x_{k+1}-x^{\star}\|^{2}
\le\|x_{k}-\tau K^{\!\top}y_{k+1}-x^{\star}\|^{2}
-2\tau\bigl(f(x_{k+1})-f(x^{\star})\bigr).
\tag{16}
\]

\noindent\textbf{Step 3.  Expand the right‑hand side of \eqref{16}.}
\[
\begin{aligned}
\|x_{k}-\tau K^{\!\top}y_{k+1}-x^{\star}\|^{2}
&=\|x_{k}-x^{\star}\|^{2}
-2\tau\inner{K^{\!\top}y_{k+1}}{x_{k}-x^{\star}}
+\tau^{2}\|K^{\!\top}y_{k+1}\|^{2}.
\end{aligned}
\tag{17}
\]

\noindent\textbf{Step 4.  Dual update expansion.}  
The dual step is $y_{k+1}= \operatorname{prox}_{\sigma g^{*}}(y_{k}+\sigma K\bar x_{k})$.
Applying Lemma~\ref{lem:prox} with $u=y_{k}+\sigma K\bar x_{k}$ and $z=y^{\star}$ gives  

\[
\inner{y_{k+1}-\bigl(y_{k}+\sigma K\bar x_{k}\bigr)}{y^{\star}-y_{k+1}}
\ge \sigma\bigl(g^{*}(y^{\star})-g^{*}(y_{k+1})\bigr).
\tag{18}
\]

After rearrangement,

\[
\|y_{k+1}-y^{\star}\|^{2}
\le\|y_{k}+\sigma K\bar x_{k}-y^{\star}\|^{2}
-2\sigma\bigl(g^{*}(y_{k+1})-g^{*}(y^{\star})\bigr).
\tag{19}
\]

\noindent\textbf{Step 5.  Expand the right‑hand side of \eqref{19}.}
\[
\begin{aligned}
\|y_{k}+\sigma K\bar x_{k}-y^{\star}\|^{2}
&=\|y_{k}-y^{\star}\|^{2}
+2\sigma\inner{K\bar x_{k}}{y_{k}-y^{\star}}
+\sigma^{2}\|K\bar x_{k}\|^{2}.
\end{aligned}
\tag{20}
\]

\noindent\textbf{Step 6.  Combine primal and dual expansions.}  
Multiply \eqref{16} by $1/(2\tau)$ and \eqref{19} by $1/(2\sigma)$, then add:

\[
\begin{aligned}
\Phi_{k+1}
&\le \Phi_{k}
-\bigl(f(x_{k+1})-f(x^{\star})\bigr)
-\bigl(g^{*}(y_{k+1})-g^{*}(y^{\star})\bigr) \\[2mm]
&\quad
-\inner{K^{\!\top}y_{k+1}}{x_{k}-x^{\star}}
+\frac{\tau}{2}\|K^{\!\top}y_{k+1}\|^{2}
+\inner{K\bar x_{k}}{y_{k}-y^{\star}}
+\frac{\sigma}{2}\|K\bar x_{k}\|^{2}.
\end{aligned}
\tag{21}
\]

\noindent\textbf{Step 7.  Re‑express the mixed terms.}  
Recall $\bar x_{k}=x_{k}+\theta (x_{k}-x_{k-1})$ and set $\theta=1$ (the
standard choice). Then  

\[
\inner{K\bar x_{k}}{y_{k}-y^{\star}}
=\inner{Kx_{k}}{y_{k}-y^{\star}}
+\inner{K(x_{k}-x_{k-1})}{y_{k}-y^{\star}}.
\tag{22}
\]

Similarly, using $y_{k+1}=y_{k}+\sigma K\bar x_{k}$ from the dual step (when $g^{*}=0$) we have  

\[
\inner{K^{\!\top}y_{k+1}}{x_{k}-x^{\star}}
=\inner{K^{\!\top}y_{k}}{x_{k}-x^{\star}}
+\sigma\inner{K^{\!\top}K\bar x_{k}}{x_{k}-x^{\star}}.
\tag{23}
\]

\noindent\textbf{Step 8.  Cancellation.}  
Insert \eqref{22}–\eqref{23} into \eqref{21}. After grouping the terms
containing $y_{k}$ and $x_{k}$ we obtain

\[
\begin{aligned}
\Phi_{k+1}
&\le \Phi_{k}
-\bigl(f(x_{k+1})-f(x^{\star})\bigr)
-\bigl(g^{*}(y_{k+1})-g^{*}(y^{\star})\bigr)\\
&\quad
-\inner{K^{\!\top}y_{k}}{x_{k}-x^{\star}}
+\inner{Kx_{k}}{y_{k}-y^{\star}}\\
&\quad
+\sigma\Bigl[
\inner{K^{\!\top}K\bar x_{k}}{x^{\star}-x_{k}}
+\inner{K\bar x_{k}}{y_{k}-y^{\star}}
+\frac{1}{2}\|K\bar x_{k}\|^{2}
\Bigr]
+\frac{\tau}{2}\|K^{\!\top}y_{k+1}\|^{2}.
\end{aligned}
\tag{24}
\]

The first two mixed inner products form the saddle‑point residual  

\[
\mathcal R_{k}:=
\inner{Kx_{k}}{y_{k}-y^{\star}}-\inner{K^{\!\top}y_{k}}{x_{k}-x^{\star}}.
\tag{25}
\]

\noindent\textbf{Step 9.  Bounding the remaining quadratic terms.}  
Using the Cauchy–Schwarz inequality and the spectral norm $\|K\|$,

\[
\begin{aligned}
\sigma\inner{K^{\!\top}K\bar x_{k}}{x^{\star}-x_{k}}
&\le \frac{\sigma\|K\|^{2}}{2}\bigl(\|\bar x_{k}\|^{2}+\|x_{k}-x^{\star}\|^{2}\bigr),\\
\sigma\inner{K\bar x_{k}}{y_{k}-y^{\star}}
&\le \frac{\sigma\|K\|^{2}}{2}\bigl(\|\bar x_{k}\|^{2}+\|y_{k}-y^{\star}\|^{2}\bigr),\\
\frac{\sigma}{2}\|K\bar x_{k}\|^{2}
&\le \frac{\sigma\|K\|^{2}}{2}\|\bar x_{k}\|^{2},\\
\frac{\tau}{2}\|K^{\!\top}y_{k+1}\|^{2}
&\le \frac{\tau\|K\|^{2}}{2}\|y_{k+1}\|^{2}.
\end{aligned}
\tag{26}
\]

Collecting the coefficients of $\|\bar x_{k}\|^{2}$, $\|x_{k}-x^{\star}\|^{2}$
and $\|y_{k}-y^{\star}\|^{2}$ we obtain a bound

\[
\Phi_{k+1}\le \Phi_{k}
-\bigl(f(x_{k+1})-f(x^{\star})\bigr)
-\bigl(g^{*}(y_{k+1})-g^{*}(y^{\star})\bigr)
+\underbrace{\bigl(\tau\sigma\|K\|^{2}\bigr)}_{\displaystyle <1}
\Bigl(\Phi_{k}+\Phi_{k+1}\Bigr).
\tag{27}
\]

\noindent\textbf{Step 10.  Rearrangement using condition \eqref{7}.}  
Since $\tau\sigma\|K\|^{2}<1$, we can move the term involving
$\Phi_{k+1}$ to the left‑hand side:

\[
(1-\tau\sigma\|K\|^{2})\Phi_{k+1}
\le (1+\tau\sigma\|K\|^{2})\Phi_{k}
-\bigl(f(x_{k+1})-f(x^{\star})\bigr)
-\bigl(g^{*}(y_{k+1})-g^{*}(y^{\star})\bigr).
\tag{28}
\]

Dividing by $1-\tau\sigma\|K\|^{2}>0$ yields

\[
\Phi_{k+1}
\le \underbrace{\frac{1+\tau\sigma\|K\|^{2}}{1-\tau\sigma\|K\|^{2}}}_{\displaystyle \alpha}
\Phi_{k}
-\frac{1}{1-\tau\sigma\|K\|^{2}}
\Bigl[
f(x_{k+1})-f(x^{\star})
+g^{*}(y_{k+1})-g^{*}(y^{\star})
\Bigr].
\tag{29}
\]

\noindent\textbf{Step 11.  Telescoping sum.}  
Iterating \eqref{29} from $k=0$ to $T-1$ and using $\alpha>1$ we obtain

\[
\sum_{k=0}^{T-1}
\Bigl[f(x_{k+1})-f(x^{\star})+g^{*}(y_{k+1})-g^{*}(y^{\star})\Bigr]
\le
\frac{1}{1-\tau\sigma\|K\|^{2}}
\bigl(\Phi_{0}-\Phi_{T}\bigr)
\le
\frac{\Phi_{0}}{1-\tau\sigma\|K\|^{2}}.
\tag{30}
\]

Dividing by $T$ and invoking convexity of $f$ and $g^{*}$ gives the
ergodic bound \eqref{9} with constant $C$ as in \eqref{10}.  This completes
the proof of Theorem~\ref{thm:main}. \hfill $\blacksquare$

\medskip
\noindent\textbf{Proof of Theorem~\ref{thm:nonerg}.}  
When $f$ is $L_{f}$‑smooth, inequality \eqref{13} applied with $u=x_{k+1}$
and $x=x_{k}$ together with the optimality condition of the proximal step
yields a descent inequality of the form  

\[
\|x_{k+1}-x^{\star}\|^{2}
\le\bigl(1-\gamma/k\bigr)\|x_{k}-x^{\star}\|^{2},
\qquad\gamma>0,
\]

and an analogous relation for $y_{k}$.  By induction we obtain the
$O(1/k)$ bound \eqref{11}.  Details follow the same algebraic steps as
above, now using the Lipschitz constant $L_{f}$ to replace the
quadratic terms in \eqref{26}.  \hfill $\blacksquare$

%%%%%%%%%%%%%%%%%%%%%%%%%%%%%%%%%%%%%%%%%%%%%%%%%%%%%%%%%%%%%%%%%%%%%%%%%%%%
\section*{Rate Derivation (With Constants)}
%%%%%%%%%%%%%%%%%%%%%%%%%%%%%%%%%%%%%%%%%%%%%%%%%%%%%%%%%%%%%%%%%%%%%%%%%%%%
From \eqref{30} we have  

\[
\mathcal G_{T}
\le\frac{C}{T},
\qquad
C=
\frac{1}{2\tau}\|x_{0}-x^{\star}\|^{2}
+\frac{1}{2\sigma}\|y_{0}-y^{\star}\|^{2}.
\]

If one chooses $\tau=\sigma=1/(2\|K\|)$ (the largest admissible equal steps),
then $C$ simplifies to  

\[
C= \|K\|\Bigl(\|x_{0}-x^{\star}\|^{2}
+\|y_{0}-y^{\star}\|^{2}\Bigr).
\]

Thus the convergence constant is linear in the norm of the operator $K$
and in the distance of the initialization from the saddle point.

%%%%%%%%%%%%%%%%%%%%%%%%%%%%%%%%%%%%%%%%%%%%%%%%%%%%%%%%%%%%%%%%%%%%%%%%%%%%
\section*{Special Cases}
%%%%%%%%%%%%%%%%%%%%%%%%%%%%%%%%%%%%%%%%%%%%%%%%%%%%%%%%%%%%%%%%%%%%%%%%%%%%
\begin{enumerate}
\item \textbf{Pure primal gradient descent.}  
If $g^{*}\equiv0$ and $K=I$, the dual update reduces to $y_{k+1}=y_{k}+ \sigma x_{k}$.
Eliminating $y$ gives the classical gradient descent with step
$\eta=\tau\sigma/(1-\tau\sigma)$, recovering the $O(1/T)$ rate of Example 1.
\item \textbf{Quadratic saddle problems.}  
For $f(x)=\frac12 x^{\top}Qx$, $g^{*}(y)=\frac12 y^{\top}Ry$ and $K$ symmetric,
the iterates can be written in closed form and the bound \eqref{9} is tight.
\item \textbf{Strongly convex $f$.}  
If $f$ is $\mu$‑strongly convex, the same analysis (with an extra
$\mu\|x_{k}-x^{\star}\|^{2}$ term) yields a linear convergence rate
$O\bigl((1-\sqrt{\mu\tau})^{k}\bigr)$.
\end{enumerate}

%%%%%%%%%%%%%%%%%%%%%%%%%%%%%%%%%%%%%%%%%%%%%%%%%%%%%%%%%%%%%%%%%%%%%%%%%%%%
\section*{Discussion}
%%%%%%%%%%%%%%%%%%%%%%%%%%%%%%%%%%%%%%%%%%%%%%%%%%%%%%%%%%%%%%%%%%%%%%%%%%%%
The PDHG algorithm enjoys a simple structure—gradient ascent on the dual
variables and a proximal descent on the primal side—yet its convergence
analysis is completely elementary: the core is the Lyapunov function
\eqref{14} and the three‑point identity \eqref{12}.  Condition \eqref{7}
guarantees that the Lyapunov function is a contraction up to the primal–dual
gap, which after telescoping yields the optimal $O(1/T)$ ergodic rate for
convex–concave saddle problems.  The proof presented above follows the
template of stochastic gradient descent (Example 1) but replaces stochastic
expectations by deterministic proximal optimality conditions, and it
explicitly tracks every algebraic manipulation with numbered steps so that
any part can be inspected independently.

\bibliographystyle{plain}
\begin{thebibliography}{9}
\bibitem{parikh2014prox}
N.~Parikh and S.~Boyd,
\newblock {Proximal Algorithms},
\newblock \emph{Foundations and Trends in Optimization}, 1(3):127–239, 2014.
\end{thebibliography}

\end{document}